\section{Energy and mesoscopic occupation}\label{sec:energy}

The energy comparison has two outputs: concentration of local particle
counts and the occupation outside functions constant on mesoscopic cells.
The latter determines the depletion except for wavelengths larger than
$R$.  The path comparisons in the following sections will control that
remaining contribution.  Throughout this section $\langle\cdot\rangle_\psi$
denotes the quadratic form in $\psi$, without division by $\|\psi\|^2$.
For a spatial cell $A$, put $q_A=n_A/(D|A|)$.

\begin{proposition}[Local energy outputs]\label{prop:local-energy}
Choose $u>0$ sufficiently small and then $0<w<u/10000$, with the
partitions of Section~\ref{sec:setup}.  In the stipulated order of
thermodynamic and dilute limits, the ground vector satisfies
\begin{equation}\label{eq:local-energy}
\begin{aligned}
 \E_{A\in\cC_R}\langle(q_A-1)^2\rangle_\Phi
       &\le CD^{-1-w/6},&
 \E_{A\in\cC_d}\langle(q_A-1)^4\rangle_\Phi
       &\le C_jD^{-1-u/6},\\
 V^{-1}\langle\dGamma(W_R)\rangle_\Phi
       &=\frac8{3\sqrt\pi}\alpha^{3/2}+o(1),&
 V^{-1}\langle\dGamma(W_d)\rangle_\Phi
       &\le\epsilon_j+o_j(1),\qquad \epsilon_j\longrightarrow0.
\end{aligned}
\end{equation}
The parameter $j$ is fixed before taking these limits.  Constants are
uniform for $\alpha$ in the fixed compact positive window.
\end{proposition}

The organization of the comparison follows the mesoscopic energy
argument of \cite[Section~3]{D}.  Its scattering-square decomposition and
Neumann cosine analysis are related to
\cite[Sections~2 and~5]{FJGMOTr}.  We give the needed estimates for the
present interaction, including the wall correction and cubic cancellation.
In particular the proof uses the bounded potential $W$ itself.

\subsection{The upper bound and Neumann boxes}

Put $c_{\rm L}=128/(15\sqrt\pi)$ and define, holding $a$ fixed when taking
derivatives,
\[
 e(y)=4\pi ay^2+4\pi c_{\rm L}a^{5/2}y^{5/2}.
\]
The Lee--Huang--Yang upper bound of
\cite[Theorem~1.1]{BCSUpper} applies because the fixed physical potential
is radial, nonnegative, compactly supported and belongs to $L^3$.
Its error is $O(\rho^{5/2+c})$ for some $c>0$, with a constant depending
on that potential.  Under our rescaling energy density is multiplied by
$\sigma^5$, and $\rho\asymp D^{-2}$.  Consequently, for some $c_*>0$,
\begin{equation}\label{eq:energy-upper}
 E_N\le V\bigl(e(D)+O(D^{-c_*})\bigr).
\end{equation}
As elsewhere, this denotes a bound on the thermodynamic limsup; no
finite-volume convergence rate is required.  To pass from the Dirichlet
upper bound to any periodic density-convergent sequence, place a
Dirichlet trial box inside the periodic box with a range margin and
side ratio tending to one.  Use a slightly larger fixed density there,
so its particle number eventually exceeds the desired one.  Slicing out
excess coordinates and averaging cannot increase the energy, since the
omitted kinetic and interaction forms are nonnegative.  Extending the
remaining slices by zero respects the periodic form domain and the range
margin makes the interactions agree.  Let the larger density decrease
to the prescribed density, using continuity of the explicit upper bound.
This proves \eqref{eq:energy-upper} along the sequences used here.

We use the same box estimates at two scales: $E=R$ yields the coarse
count variance, and $E=B$ yields the microcell fourth moment and both
occupation estimates.  For a generic box side $E$, abbreviate $Q=E$
and $v_b=E^3$.
Thus $Q\asymp D^\theta$ for $\theta=u$ or $w$.  Restrict configuration
space to sectors assigning each label to a box and discard interactions
between boxes.  These are Neumann form restrictions: we restrict gradient
integrals, without differentiating the sector indicators.  Each local
vector is symmetric in the labels assigned to its box, and may take
values in the Hilbert space of the spectator coordinates.

It is convenient to reflect the interaction at each wall.  For a radial
function $g$ supported in $\{|x|\le r\}$ define on $[0,E]^3$
\[
 g^s(x,y)=\sum_{z\in\Z^3}g(p_z(x)-y),
\]
where, coordinatewise, $p_z$ maps $[0,E]^3$ onto the cube translated by
$Ez$, reversing a coordinate precisely when the corresponding $z_i$ is
odd.  Write $H_b$ for the Neumann Hamiltonian with pair potential $W^s$.
The added mirror terms are nonnegative, but their expectation on the
original ground vector is small:
\begin{equation}\label{eq:energy-mirror-error}
 \sum_b\langle H_b\rangle\le E_N+CV/E.
\end{equation}
Indeed an added interaction requires two particles within $Cr$ of each
other and within $Cr$ of a wall.  Its multiplicity is bounded.  At most
$C(V/E)r^{-2}$ cubes of side comparable with $r$ cover all walls.  The
small-cube pair estimate proved before \eqref{eq:pair-bound} gives
$CD^2r^6$ ordered pairs per such cube, and $\|W\|_\infty\le CD^2$.
Since $r\asymp D^{-1}$, the resulting bound is $CV/E$.
All later excitation localizations use this same mirror Hamiltonian.

Let $u_p$, $p\in(\pi/E)\N_0^3$, be the real normalized Neumann cosines:
\[
 u_p(x)=v_b^{-1/2}\prod_{i=1}^3c_{p_i}\cos(p_ix_i),\qquad
 c_0=1,\quad c_s=\sqrt2\ (s>0).
\]
Even extension shows that the integral operator with kernel $g^s$ has
eigenvalue $\widehat g(p)$ on $u_p$, with Fourier convention
$\widehat g(p)=\int g(x)e^{-ip\cdot x}\,\dd x$.  In particular its row
integral is the constant $\widehat g(0)$.
Write $a_p,a_p^*$ for the corresponding bosonic operators and set
\begin{equation}\label{eq:energy-bands}
 J=Q^{100},\qquad
 n_+=\sum_{p\ne0}a_p^*a_p,\qquad
 n_L=\sum_{0<|p|\le J/E}a_p^*a_p,\qquad
 n_H=n_+-n_L,\qquad M\asymp Dv_bQ^{-400}.
\end{equation}
Take $M$ integral.  We refer to $0<|p|\le J/E$ as the low band and the
remaining nonzero modes as the high band.  All power inequalities below
hold if $u$ is sufficiently small, for example $u<10^{-9}$ and
sufficiently small relative to $c_*$.  Further decreases are harmless.

\subsection{A rough bound retaining count and kinetic penalties}

We first obtain a lower bound with an error much larger than the
Lee--Huang--Yang term.  Its purpose is to control vectors with many
excitations, leaving only a small-excitation sector for the sharper
calculation.  The construction is the radial replacement and small-box
method of Dyson and Lieb--Yngvason~\cite{Dyson,LiebYngvason}, with an
explicit reserve of kinetic energy.

\begin{lemma}[Rough box bound]\label{lem:energy-rough}
For $0\le n\le6Dv_b$, as a quadratic-form inequality on the bosonic
$n$-particle Neumann space,
\begin{equation}\label{eq:rough-energy}
\begin{split}
 H_b\ge{}&\frac{4\pi an^2}{v_b}-Cv_bD^{1-0.001}
       +c\,\dGamma(\min\{p^2,1\})\\
 &+cD\sum_{A\in\cC_d,\ A\subset b}|A|
    \mathfrak h\left(q_A-\frac{n}{Dv_b}\right),
 \qquad \mathfrak h(y)=\frac{y^2}{1+|y|}.
\end{split}
\end{equation}
The same inequality holds with Hilbert-space-valued spectator variables.
\end{lemma}

\begin{proof}
We can discard the mirror terms for this estimate.  Subdivide each
$d$-cell into equal cubes of side $\ell\asymp D^{-0.31}$, and put
$b_1=D^{-0.36}$ and $\xi=D^{-0.002}$.  For large $D$, $r<b_1/2$.
Restrict to the corresponding assignment sectors and drop interactions
between small cubes.

Split the kinetic coefficient as $1=(1-2\xi)+\xi+\xi$.
The radial replacement below uses the first term only within distance
$b_1$ of another particle in the small cube.  We use one $\xi$ reserve
for Temple's inequality and retain the other for the kinetic term in
\eqref{eq:rough-energy}; outside those neighborhoods the first term
also remains available.

Fix one particle coordinate and all the others in its small cube.
Partition the cube into the Voronoi regions of those other particles,
clipped at the walls.  Each region is star-shaped about its center.
The radial scattering variational principle on a ray gives, for
$R'\ge r$,
\begin{equation}\label{eq:energy-ray}
 \int_0^{R'}\bigl((1-2\xi)|f'(s)|^2+\tfrac12W(s)|f(s)|^2\bigr)
                  s^2\,\dd s
 \ge(1-2\xi)a|f(R')|^2.
\end{equation}
With kinetic coefficient one the exact endpoint minimum is
$a(1-a/R')^{-1}|f(R')|^2$: integrate the scattering equation against its
regular solution, then complete the quadratic form around the minimizer.
To obtain \eqref{eq:energy-ray}, first bound the potential term below by
$(1-2\xi)W/2$, and then discard $(1-a/R')^{-1}\ge1$.
No monotonicity of the radial function $W$ is involved.

Let $U$ be constant on $b_1/2<|x|<b_1$, zero elsewhere, and normalized
by $\int U=4\pi$.  Average \eqref{eq:energy-ray} over radii available
on each ray with weight $U(R')(R')^2\,\dd R'$.  The weight has total at
most one.  Integrating the rays replaces the short-distance energy by
$(1-2\xi)aU(d_{\rm nn})$, where $d_{\rm nn}$ is the nearest-neighbor
distance in the small cube; the expression is zero when there is no
neighbor.  Assigning half of each pair potential to each endpoint uses
no interaction twice.  This argument first holds for smooth vectors
and then by form approximation.

In a small cube with $k\le C_0D\ell^3$ particles, the mean of the
replacement sum in its normalized constant vector is at least
\[
 \frac{4\pi(1-2\xi)ak(k-1)}{\ell^3}
       \left(1-C\frac{b_1}{\ell}-C\frac{kb_1^3}{\ell^3}\right).
\]
For this estimate put one coordinate at least $b_1$ from the walls,
integrate a possible nearest neighbor over the shell, and exclude closer
neighbors by the union bound.  Distinct choices of the nearest neighbor
are disjoint outside null sets.  The replacement sum is nonnegative
and bounded by $Ckab_1^{-3}\le CD^{0.15}$.  Its variance in the constant
vector is therefore at most this upper bound times its mean.  The
kinetic gap supplied to Temple's inequality is
$\pi^2\xi/\ell^2\asymp D^{0.618}$.
Temple's inequality yields the lower bound
\[
 \frac{4\pi a}{\ell^3}(1-CD^{-0.002})k(k-1).
\]
Here the other relative errors are
$b_1/\ell=O(D^{-0.05})$ and $Db_1^3=D^{-0.08}$; the variance correction
is smaller still.  The formula also holds for $k=0,1$.

For $k>C_0D\ell^3$, partition the labels into groups of sizes between
$C_0D\ell^3/3$ and $C_0D\ell^3$, allowing integer rounding, and drop
intergroup interactions.  Choose a fixed cutoff $b_*>12$, and then $C_0$
sufficiently large.  With
\[
 F(y)=\begin{cases}y^2,&0\le y\le b_*,\\2b_*y-b_*^2,&y>b_*,\end{cases}
\]
the same estimate in every number sector is bounded below by
\[
 \frac{4\pi a}{\ell^3}(1-CD^{-0.002})
       \left[(D\ell^3)^2F\left(\frac{k}{D\ell^3}\right)-k\right].
\]
Indeed $F(y)\le y^2$ below the group threshold, whereas above it the
group estimate has a linear coefficient larger than $2b_*D\ell^3$.
For $\vartheta=n/(Dv_b)\in[0,6]$,
\[
 F(y)\ge\vartheta^2+2\vartheta(y-\vartheta)
                         +c\mathfrak h(y-\vartheta),\qquad y\ge0.
\]
Below $b_*$ the remainder is a square.  Above $b_*$, writing
$s=y-\vartheta$ and $a_0=b_*-\vartheta>0$, it is
$2a_0s-a_0^2\ge a_0s$.  This proves the displayed comparison uniformly.
Summing small cubes cancels the tangent deviations.  The self-count
terms cost at most $Ca\ell^{-3}n\le Cv_bD^{0.93}$, and the relative
loss in the leading scalar costs $Cv_bD^{1-0.002}$.
Convexity of $\mathfrak h$ and Jensen's inequality inside each $d$-cell
now give the scalar and count terms of \eqref{eq:rough-energy}.

It remains to recover a fixed amount of the indicated kinetic form.
The preceding construction leaves coefficient $\xi$ everywhere and
$1-2\xi$ outside the $b_1$-neighborhoods of the other particles.
For a coordinate slice enlarge those neighborhoods to their union $S$;
then $|S|\le Cnb_1^3$.  The vector fields $\nabla u_p/|p|$ are orthonormal.
Their projection $\Pi$ for $0<|p|\le D^{0.003}$ has squared evaluation
norm at most $CD^{0.009}$.  Hence, by duality,
\[
 \|\Pi\1_S F\|_2^2\le CD^{0.009}nb_1^3\|F\|_2^2
                \le CD^{-0.071}Q^3\|F\|_2^2=o(\xi)\|F\|_2^2.
\]
Apply this to $F=\nabla f$ and use
$\|\Pi F\|_2^2\le2\|\1_{S^c}F\|_2^2+2\|\Pi\1_SF\|_2^2$.
Neumann integration by parts identifies the projected gradient norm
with $\sum_{0<|p|\le D^{0.003}}p^2|\langle u_p,f\rangle|^2$.
Above this cutoff the everywhere reserve suffices because
$\xi p^2\ge D^{0.004}$.  Sum over coordinates.  This last argument is
on the original Neumann form domain of the $E$-box, after summing back
the small-cube restrictions; it does not impose artificial Neumann
conditions on an arbitrary collection of small-cube functions.
\end{proof}

\subsection{Restricting the excitation number}

The cutoff below is a finite-band instance of the matrix localization
method of \cite[Appendix~A]{LiebSolovej}.  We give its form estimate
directly so that the retained positive terms remain explicit.

Comparing Lemma~\ref{lem:energy-rough} with the tangent to
$4\pi ay^2$ at $D$, and using \eqref{eq:energy-upper} and
\eqref{eq:energy-mirror-error}, gives the preliminary estimate
\begin{equation}\label{eq:energy-preliminary}
 V^{-1}\sum_{b,\ n\le6Dv_b}\langle n_+\rangle\le CQ^2D^{1-0.001}.
\end{equation}
For completeness the higher number sectors do not spoil this comparison:
partition their labels into groups of sizes between $2.5Dv_b$ and
$6Dv_b$, apply the rough bound to each group and drop intergroup
interactions.  The energy per particle is then at least
$10\pi\alpha-o(1)$, exceeding the tangent slope $8\pi\alpha$ by a
fixed positive amount.  Thus those sectors have a surplus of order $n$.
For $n\le6Dv_b$, the scalar tangent remainder is nonnegative and
$\min(p^2,1)\ge cE^{-2}$ on the nonconstant modes, proving
\eqref{eq:energy-preliminary}.

Localize only the middle number sectors
\[
                       Dv_b/2\le n\le6Dv_b.
\]
Choose real Lipschitz functions $f,h$ with $f^2+h^2=1$, $f=1$ below
$M/8$, $f=0$ above $M/4$, and Lipschitz constants at most $C/M$.
The three localizing operators are
\begin{equation}\label{eq:energy-cutoffs}
 F_1=f(n_L)^2,\qquad F_2=h(n_L)^2,\qquad F_3=\sqrt2f(n_L)h(n_L).
\end{equation}
They have $\sum F_i^2=1$.  The first and third pieces are supported on
$n_L\le M/4$ and will be called low pieces; the second is supported on
$n_L\ge M/8$.  Also $\psi=F_1\psi+F_2\psi$, which later transfers
count estimates even though counts need not commute with these cutoffs.
Finite-rank smooth one-body projections preserve the Neumann form domain.

Here is the localization error, with $V_n^s$ the sum of mirror pair
potentials:
\begin{equation}\label{eq:energy-cutoff-error}
 \left|\sum_{i=1}^3\langle H_b\rangle_{F_i\psi}
                      -\langle H_b\rangle_\psi\right|
 \le\frac C{M^2}\left(\langle V_n^s\rangle_\psi
       +\|W\|_1\frac{J^3}{v_b}n\langle n_+\rangle_\psi\right).
\end{equation}
To prove it, let $\Pi_k$ project onto $n_L=k$.  Kinetic energy commutes
with the cutoffs, and the potential has block width two.  The difference
of its forms is
\[
 -\frac12\sum_{k,l,i}(F_i(k)-F_i(l))^2
                       \langle\Pi_k\psi,V_n^s\Pi_l\psi\rangle.
\]
Positivity, form Cauchy--Schwarz and the finite block width bound this
by $CM^{-2}$ times the block-diagonal expectation.  The latter is the
phase average of $e^{-isn_L}V_n^se^{isn_L}$.  For a pair only the two
coordinate rotations matter.  Expanding each as
$1+(e^{is}-1)P_L$, every additional term has a low projection on at
least one coordinate.  The evaluation bound
$P_L(x,x)\le CJ^3/v_b$ and the constant row integral of $W^s$ give
\[
 P_{L,i}W^s(x_i,x_k)P_{L,i}
       \le C\|W\|_1J^3v_b^{-1}P_{L,i}.
\]
Expanding squared norms and summing pairs proves
\eqref{eq:energy-cutoff-error}.
The analogous error for $\dGamma(A)$, $\|A\|\le1$, is at most
$CnM^{-2}\|\psi\|^2$, by block width one and its norm bound $n$.

The summed potential expectation is $O(VD)$, including the mirror
correction.  Together with \eqref{eq:energy-preliminary},
\eqref{eq:energy-cutoff-error} has total error $O(VQ^{-1})$.
Even replacing $\|W\|_1\le C/D$ by the weaker $CQ$ gives the per-volume
bound
\[
 O\bigl(D^{-1}Q^{794}+D^{-0.001}Q^{1097}\bigr)=O(Q^{-1}).
\]
All these inequalities include sector weights and spectator integration.

\subsection{The sharp box estimate}

We now analyze a low piece.  The aim is to recover the
Lee--Huang--Yang scalar while retaining two positive terms: a weak
particle-excitation gap and a quasiparticle energy.  The latter will
identify the occupation observable.

Let $\varphi$ be the scattering function of $W$, and set
\[
 \omega=1-\varphi,\qquad g=W\varphi,\qquad m=g\omega.
\]
The maximum principle and the radial scattering equation give
$0\le\varphi\le1$, $\widehat g(0)=8\pi a$ and $\|m\|_1\le C/D$.
For $p\ne0$ put
\begin{equation}\label{eq:energy-tau}
 S=E^{-2}(n_++Jn_H),\qquad
 \tau_p=p^2-E^{-2}\bigl(1+J\1_{\{p\ \text{high}\}}\bigr).
\end{equation}
Thus $\tau_p\asymp p^2$.  For the constant mode use coherent vectors
$|z\rangle=e^{-|z|^2/2}\sum_{q\ge0}z^q|q\rangle/\sqrt{q!}$,
measure $\dd\lambda(z)=\pi^{-1}\dd^2z$, and coherent projections
$\psi_z=\langle z|\psi\rangle$ onto the excitation Fock space.
For $|z|^2\le2n$ define
\[
 \rho_z=|z|^2/v_b,\qquad t_p=\rho_z\widehat g(p),\qquad
 D_p=(\tau_p^2+2t_p\tau_p)^{1/2}.
\]
These roots are positive: $|t_p|\le C$, while $\widehat g(p)\ge0$ on
every fixed momentum ball for large $D$, by positivity and shrinking
support.  Large $p$ are controlled by $\tau_p\asymp p^2$.
After rotating excitation modes by the phase of $z$, the change of
variables used below diagonalizes the one-mode quadratic form
\[
 (\tau_p+t_p)a_p^*a_p+\tfrac12t_p(a_p^2+a_p^{*2}).
\]
Specifically, define $b_p$ through
\begin{equation}\label{eq:energy-bogoliubov}
\begin{gathered}
 a_p=U_pb_p+V_pb_p^*,\qquad U_p>0,\quad U_p^2=1+V_p^2,\\
 V_p^2=\tfrac12\left(\frac{\tau_p+t_p}{D_p}-1\right),\qquad
 2U_pV_p=-t_p/D_p.
\end{gathered}
\end{equation}
With this choice,
\[
 (\tau_p+t_p)a_p^*a_p+\tfrac12t_p(a_p^2+a_p^{*2})
  =D_pb_p^*b_p+\tfrac12(D_p-\tau_p-t_p).
\]
The formulas with $z=0$ have $U_p=1,V_p=0$.

\begin{lemma}[Sharp energy with positive remainders]\label{lem:energy-sharp}
If $Dv_b/2\le n\le6Dv_b$ and $\psi$ is supported on $n_L\le M/4$,
then
\begin{equation}\label{eq:sharp-energy}
\begin{split}
 \langle H_b\rangle_\psi\ge{}&
 \bigl(v_be(n/v_b)-Cv_bQ^{-1/4}\bigr)\|\psi\|^2+c\langle S\rangle_\psi\\
 &+c\int_{|z|^2\le2n}\sum_{p\ne0}\widetilde D_p
                       \|b_p\psi_z\|^2\,\dd\lambda(z),\\
 &\hspace{12mm}\widetilde D_p=D_p\ \text{on the low band},\qquad
 \widetilde D_p=D_p/J\ \text{on the high band}.
\end{split}
\end{equation}
\end{lemma}

The proof has three parts.  Scattering-square splitting produces a
quadratic Hamiltonian and a cubic term.  Its quadratic vacuum energy is
the Lee--Huang--Yang scalar.  Completing the high-frequency cubic square
creates a contraction that cancels against the retained $m$ term.
We supply each part, since the cancellation must leave the positive
remainders in \eqref{eq:sharp-energy}.

\begin{proof}[Proof of Lemma~\ref{lem:energy-sharp}]
\emph{Scattering-square splitting.}
Let $P=|u_0\rangle\langle u_0|$, $Z=P_i^\perp P_j^\perp$ and $A=1-Z$.
For every summand in the mirror kernel, with its argument suppressed,
there is the exact identity of forms
\begin{equation}\label{eq:energy-scattering-square}
 W=(Z+\omega A)^*W(Z+\omega A)+ZgA+AgZ+A(g+m)A.
\end{equation}
It follows by expanding, using $g=W(1-\omega)$ and
$m=W\omega(1-\omega)$.
In the cubic blocks $P_iP_j^\perp g^sZ$ and their adjoints retain only
the low projection in coordinate $j$.  The high remainder costs a small
fixed fraction of the positive squares and $Cn_H$.  Here is a direct
verification.  In each mirror term substitute
$Z=(Z+\omega A)-\omega A$.  Cauchy--Schwarz against the square bounds
the first part by that square plus $Cn_H$: the latter bound uses
$g^2/W\le W$, the constant row integral of $W^s$, and
$n\|W\|_1/v_b\le C$.  The other part contains $m^sA$.
Between $P_iP_j^{\rm high}$ and
$A=P_iP_j+P_iP_j^\perp+P_i^\perp P_j$ its first block vanishes.
Its other blocks act only on high modes, by the constant row integral
for the direct block and cosine diagonalization for the exchange block.
Their norms are at most $\|m\|_1/v_b$, again giving $Cn_H$ after summing.
This is absorbed by $E^{-2}Jn_H$ since $J/E^2\to\infty$.

The retained $g$ direct terms in the number-$n$ sector combine into
\begin{equation}\label{eq:energy-number-term}
 \frac{\widehat g(0)}{2v_b}
            \{n(n-1)-n_+(n_+-1)\}.
\end{equation}
After separating $S$, the remaining terms are
\begin{equation}\label{eq:energy-retained-hamiltonian}
\begin{split}
 &\frac{\widehat m(0)}{2v_b}a_0^{*2}a_0^2\\
 &+\sum_{p\ne0}\left[
  \tau_pa_p^*a_p+
  \frac{\widehat g(p)+\widehat m(p)+\widehat m(0)}{v_b}
                           a_0^*a_p^*a_pa_0
  +\frac{\widehat g(p)}{2v_b}(a_0^{*2}a_p^2+a_p^{*2}a_0^2)\right]\\
 &+\frac1{v_b}\sum_{\substack{p\ \text{low}\\ k,l\ne0}}
 C_{pkl}\widehat g(k)
                (a_0^*a_p^*a_la_k+a_k^*a_l^*a_pa_0),
 \qquad C_{pkl}=\sqrt{v_b}\int_bu_pu_ku_l.
\end{split}
\end{equation}
These are the matrix elements of
\eqref{eq:energy-scattering-square}: the constant mode paired through
$g^s$ with mode $k$ leaves exactly the triple cosine coefficient in the
last line.  Blocks with just one nonconstant index vanish.  The
coefficients $C_{pkl}$ are bounded, and fixing any two indices leaves at
most a fixed number of possible third indices.  These facts follow
coordinatewise from the product-to-sum identity for cosines.  We keep
these exact coefficients, including modes with zero components.

On the present excitation cutoff,
\begin{equation}\label{eq:energy-number-square}
 n_+^2\le C(Mn_++nn_H).
\end{equation}
Thus replacing \eqref{eq:energy-number-term} by $4\pi an^2/v_b$ costs
$O(1)\|\psi\|^2$ and an $o(1)$ fraction of $S$.
For example its two excitation coefficients are $O(Q^{-400})$ on
$n_+$ and $O(1)$ on $n_H$, compared with $E^{-2}$ and $JE^{-2}$.

\emph{Coherent resolution.}
Use the resolution of the identity in the zero mode.  The upper symbols
of $a_0^*a_0$ and $a_0^{*2}a_0^2$ are respectively
$|z|^2-1$ and $|z|^4-4|z|^2+2$; these follow by scalar gamma integration,
or see \cite{LSY2005}.  Hence in
\eqref{eq:energy-retained-hamiltonian} we may substitute $z,\bar z$
for $a_0,a_0^*$, with total integrated error $O(1)\|\psi\|^2$.
Indeed the Fourier coefficients are $O(D^{-1})$, $n\le6Dv_b$, and
\[
 \int|z|^2\|\psi_z\|^2\,\dd\lambda(z)
       =\langle n-n_++1\rangle_\psi.
\]
The region $|z|^2>2n$ is negligible, also with fixed polynomial weights:
after angular integration its radial variables are gamma variables
with integer shapes at most $n+1$, whose tails there are exponentially
small.  All negative terms, estimated without kinetic energy, have
polynomial operator bounds on sectors with at most $n$ excitations.
For the pairing terms this follows from Parseval for the bounded mirror
kernel; there are only polynomially many low cubic indices.
The coherent projections preserve both $n_L\le M/4$ and $n_+\le n$.
It therefore suffices to estimate a normalized excitation vector $\eta$
in that subspace, at $|z|^2\le2n$.  For this calculation rotate phases
so that $z$ is real.

\emph{The quadratic vacuum energy.}
Diagonalizing the $g$ quadratic terms gives positive
$\sum D_p\|b_p\eta\|^2$, the remaining $m$ terms, and the scalar
\begin{equation}\label{eq:energy-lhy-scalar}
 \frac12\sum_{p\ne0}(D_p-\tau_p-t_p)
       +\frac12v_b\rho_z^2\widehat m(0)
 =4\pi c_{\rm L}a^{5/2}v_b\rho_z^{5/2}+O(v_bQ^{-1/4}).
\end{equation}
We justify the error without any regularity assumption on $W$.
The equation $-\Delta\omega=g/2$ and the Newton kernel give
\[
 \widehat m(0)=\frac12\int_{\R^3}\frac{\widehat g(k)^2}{|k|^2}
                      \,\dd\bar k,\qquad
                   \dd\bar k=(2\pi)^{-3}\dd k.
\]
Writing $H(k)=\widehat g(k)/(8\pi a)$, the radial sine formula,
$\|g\|_\infty\le CD^2$ and $\supp g\subset\{|x|\le C/D\}$ imply
\begin{equation}\label{eq:energy-fourier-decay}
 |H(k)|\le\min\{1,CD/|k|\},\qquad
 |\nabla H(k)|\le C\min\{D^{-1},|k|^{-1}\}.
\end{equation}
The uniform derivative bound uses the first moment of $g$.
For the large-frequency bounds differentiate
$4\pi\int_0^r g(s)s\sin(|k|s)/|k|\,\dd s$ and use the height and range
bounds; no derivative of $g$ is taken.

Neumann lattice summation in the positive octant gives
\begin{equation}\label{eq:energy-riemann-m}
 \frac1{v_b}\sum_{p\ne0}\frac{H(p)^2}{p^2}
   =\int_{\R^3}\frac{H(p)^2}{p^2}\,\dd\bar p
                    +O\left(\frac{1+\log D}{E}\right).
\end{equation}
To see this, compare cubes of mesh $\pi/E$.  The octant weight
$\pi^{-3}$ equals the full-space weight $(2\pi)^{-3}$ for the even
integrand.  The cubes meeting zero cost $O(E^{-1})$.  Away from zero,
the derivative is bounded by
$C(|p|^{-3}+D^{-1}|p|^{-2})$ up to radius $D$ and by
$C(D^2|p|^{-5}+D|p|^{-4})$ above it.  Shell summation, including face
cubes, proves \eqref{eq:energy-riemann-m}; here $\log E=O(\log D)$.
Thus the second term in \eqref{eq:energy-lhy-scalar} may be replaced
by $\sum t_p^2/(4p^2)$ at the stated accuracy.

Replacing $\tau_p$ by $p^2$ in the first term costs $O(v_b/E)$.
At bounded frequencies the derivative of the summand is bounded by
$C(1+|p|^{-1})$ and the shift is $E^{-2}$.  At large frequencies the
derivative is $O(|p|^{-4})$; the additional high-band shift $J/E^2$
is compensated by summing only above $J/E$.
In the combined summand replace $t_p$ by $t_0=8\pi a\rho_z$.
Above radius $Q$ its absolute tail is $O(Q^{-1})$ per volume by Taylor
expansion.  Below $Q$, positivity and support give
$|t_p-t_0|\le Cp^2/D^2$, while differentiation in $t$ costs at most
$C(1+p^{-2})$.  The summed error is $O(Q^5/D^2)$ per volume.
Finally, the constant-$t_0$ Riemann sum differs from its integral by
$O((1+\log E)/E)$ per volume: near zero use the bounds
$C|p|^{-2},C|p|^{-3}$ for the function and its derivative, and at infinity
use the Taylor bounds $C|p|^{-4},C|p|^{-5}$.
All these errors are $O(Q^{-1/4})$.
The integral equals
\[
 \frac{t_0^{5/2}}{4\pi^2}
 \int_0^\infty x^2
   \left(\sqrt{x^4+2x^2}-x^2-1+\frac1{2x^2}\right)\dd x.
\]
The radial integral is $8\sqrt2/15$, by evaluating an elementary
antiderivative and taking its endpoint limits.  Substitution of
$t_0=8\pi a\rho_z$ proves \eqref{eq:energy-lhy-scalar}, including $z=0$.

The scalar is now correct.  We next control the cubic term while
preserving a fraction of the quasiparticle energy.  The $m$ terms still
present in \eqref{eq:energy-retained-hamiltonian} are essential here.

\emph{The cubic terms with low $k$.}
Put $\lambda=(Dv_b)^{-1/2}$ and $N_+=\langle n_+\rangle_\eta$.
After coherent substitution the cubic coefficients are bounded by
$C\lambda|H(k)|$.  Whenever we sum triples below, $p$ is low and
$C_{pkl}\ne0$.
For low $k$, Cauchy--Schwarz applied to
$\|a_l^*a_p\eta\|$ and $\|a_k\eta\|$ gives
\begin{equation}\label{eq:energy-low-cubic}
 C\lambda J^{3/2}(M+CJ^3)^{1/2}N_+=o(E^{-2})N_+.
\end{equation}
Indeed there are $O(J^3)$ low indices, any two indices have bounded
multiplicity for the third, and
$\|a_l^*a_p\eta\|^2=\|a_la_p\eta\|^2+\|a_p\eta\|^2$.
The double-annihilation sum is bounded by
$\langle n_Ln_+\rangle_\eta$.

\emph{The high $k$ creation part.}
For high $k$ use $a_k=U_kb_k+V_kb_k^*$, with
$|U_k|\le C$ and $|V_k|\le C|H(k)|/k^2$.
Commute $b_k^*$ to the left in the $V_k$ part.  Besides the reordered
term this produces the commutator with $a_l$ when $l=k$; the
commutator with $a_p^*$ vanishes because $p$ is low.  Set
\[
 T_H=\sum_{k\ \text{high}}\frac{k^2}{J}\|b_k\eta\|^2.
\]
There are $O(J^3)$ triples for fixed $k$.  Weighted Cauchy--Schwarz
therefore bounds the reordered term by
\[
 C\lambda J^2 T_H^{1/2}
       \left(\sum_{p,k,l}\frac{H(k)^4}{k^6}
                                \|a_p^*a_l\eta\|^2\right)^{1/2}.
\]
The unweighted squared-norm sum is at most
$C(M+J^3)N_+$, again by cosine multiplicities and commutation.
Young's inequality bounds this by an arbitrarily small fixed fraction
of $T_H$ plus
\begin{equation}\label{eq:energy-high-creation-error}
 C\lambda^2J^4(E/J)^6(M+J^3)N_+=o(E^{-2})N_+.
\end{equation}
For the commutator, $C_{pkk}\ne0$ requires some nonzero component of
$p$ to equal twice the corresponding component of $k$.
For fixed $p$, the sum of $H(k)^2/k^2$ over these high $k$ is at most
$CE^2(1+\log D)$: on each such plane use dyadic shells and the decay
in \eqref{eq:energy-fourier-decay} above $D$.
Its total contribution is at most
\[
 C\lambda E^2(1+\log D)J^{3/2}\sqrt{N_+}.
\]
It costs an arbitrarily small fixed fraction of $E^{-2}N_+$ and the
scalar error $CD^{-1}Q^{303}(1+\log D)^2$.

\emph{The high $k$ annihilation part and its contraction.}
Complete a square between the $U_kb_k$ cubic and the fraction $1-3/J$
of the high diagonal energy.  The resulting negative form is bounded by
\begin{equation}\label{eq:energy-cubic-square}
 \frac{\rho_z}{v_b}\sum_{k\ \text{high}}
  \frac{\widehat g(k)^2U_k^2}{(1-3/J)D_k}
       \left\|\sum_{p,l}C_{pkl}a_l^*a_p\eta\right\|^2.
\end{equation}
Normal ordering separates its quartic and contraction parts.
The absolute value of the quartic part is at most
\begin{equation}\label{eq:energy-cubic-quartic}
 C\lambda^2(E/J)^2J^3\langle n_Ln_+\rangle_\eta
                                      =o(E^{-2})N_+.
\end{equation}
For fixed $p,p'$ the sums over $k,l,l'$ in the products
$\|a_{l'}a_p\eta\|\,\|a_la_{p'}\eta\|$ have bounded multiplicity.
Put $A_p=\sum_{l\ne0}\|a_la_p\eta\|^2$.  Cauchy--Schwarz first
bounds each fixed-$p,p'$ sum by $C\sqrt{A_pA_{p'}}$, and then
\[
 \sum_{p,p'\ {\rm low}}\sqrt{A_pA_{p'}}
 \le \#\{p:p\ {\rm low}\}\sum_{p\ {\rm low}}A_p
 \le CJ^3\langle n_Ln_+\rangle_\eta .
\]
The weight is bounded using $|k|>J/E$.  This proves
\eqref{eq:energy-cubic-quartic}.

The contraction matrix for fixed $k$ has entries
\[
 B^{(k)}_{pp'}=\sum_{l\ne0}C_{pkl}C_{p'kl}
\]
on low modes.  It is positive with bounded norm, since it is the
compression of multiplication by $v_bu_k^2$, whose supremum is at
most eight.  The omitted $l=0$ term is zero when $k$ is high and $p$
is low.  If every component $k_i>2J/E$, the product-to-sum formula for
$u_k^2$ makes this compression exactly the identity.
For all high $k$,
\[
 \frac{U_k^2k^2}{(1-3/J)D_k}\le1+C/J.
\]
The exceptional component slabs satisfy
\begin{equation}\label{eq:energy-slab-sum}
 \frac1{v_b}\sum_{\substack{k\ \text{high}\\\min_i k_i\le2J/E}}
                \frac{H(k)^2}{k^2}
       \le C\frac{J(1+\log D)}E.
\end{equation}
There are $O(J)$ lattice planes per slab; each plane has the same
$CE^2(1+\log D)$ bound used for the commutator.  Combining
\eqref{eq:energy-riemann-m} and \eqref{eq:energy-slab-sum}, the
contraction in \eqref{eq:energy-cubic-square} is at most
\begin{equation}\label{eq:energy-cubic-contraction}
 \left(2\rho_z\widehat m(0)
       +C\left[J^{-1}+D^{-1}J(1+\log D)/E\right]\right)
                                      \langle n_L\rangle_\eta.
\end{equation}
The factor two comes from
$\int\widehat g(k)^2/k^2\,\dd\bar k=2\widehat m(0)$.
On the other hand the low-band $m$ diagonal in
\eqref{eq:energy-retained-hamiltonian}, after coherent substitution,
is at least
\[
 2\rho_z\widehat m(0)n_L-CD^{-2}(J/E)^2n_L.
\]
This uses the second-moment estimate from $\supp m\subset\{|x|\le r\}$.
Its high-band part is nonnegative because
$\widehat m(p)+\widehat m(0)\ge0$ for the nonnegative function $m$.
Thus the leading term in \eqref{eq:energy-cubic-contraction} cancels,
and every remaining coefficient is $o(E^{-2})$.

For clarity, the preceding absorptions have ample strict margins.
The low-$k$ coefficient in \eqref{eq:energy-low-cubic} is bounded by
$C(Q^{-50}+D^{-1/2}Q^{298.5})$; the coefficient in
\eqref{eq:energy-high-creation-error} is
$C(Q^{-594}+D^{-1}Q^{103})$; and
\eqref{eq:energy-cubic-quartic} has coefficient $CQ^{-298}$.
The commutator scalar is $o(v_bQ^{-1/4})$ for sufficiently small $u$.
Completing the main square used only $1-3/J$ of the high diagonal;
the creation estimate uses a small fraction of the remaining
$J^{-1}$ energy.  Thus a fixed multiple of $\widetilde D_p$ remains,
as does a fixed fraction of $S$.

\emph{Integration and domains.}
Integrate the estimates on coherent vectors.  The function
$4\pi c_{\rm L}a^{5/2}\rho^{5/2}$ is convex, and its derivative at
$\rho=n/v_b$ is $O(D^{-1})$.  Its tangent there, together with the
coherent second-moment identity, shows that replacing its integrated
value by the value at $n/v_b$ loses at most
$CD^{-1}(\|\psi\|^2+\langle n_+\rangle_\psi)$.
This is absorbed in the scalar error and the retained $S$.
The gamma tail previously discarded is exponentially small.

All calculations can first be made with finitely many excitation modes.
At fixed parameters the bounded interaction is form-continuous under
finite-mode approximation, which preserves the excitation cutoffs.
The pairing coefficients have square-summable Fourier coefficients by
Parseval, and the cubic has finitely many low indices.  The displayed
bounds justify passage in the completed squares; a positive remainder
can first be retained on finitely many modes and then passed to the
limit by lower semicontinuity.  This proves
\eqref{eq:sharp-energy} on the full form domain.
\end{proof}

\subsection{The energy surplus and count moments}

The sharp estimate now applies to the low pieces; the rough estimate
handles their complement and the outlying particle-number sectors.
The bookkeeping is most transparent relative to the affine function
\[
                  t_b(n)=v_be(D)+e'(D)(n-Dv_b).
\]
Its sum over boxes is exactly $Ve(D)$ in every assignment sector.

\begin{lemma}[Positive quantities paid for by the energy surplus]
\label{lem:energy-surplus}
After summing boxes, assignment sectors, and spectator variables, each
of the following quantities is bounded by $CVQ^{-1/5}$:
\begin{enumerate}
\item the original box-count penalty
$\sum_bDv_b\langle\mathfrak h(n/(Dv_b)-1)\rangle$;
\item on high pieces, the sum of
$\langle\dGamma(\min\{p^2,1\})\rangle$ and
$D\sum_{A\subset b}|A|\langle\mathfrak h(q_A-1)\rangle$;
\item on low pieces, the $S$ form and the integrated positive
quasiparticle form in \eqref{eq:sharp-energy};
\item on number sectors outside $[Dv_b/2,6Dv_b]$, the sum of
$(n+Dv_b)$ times their squared norms.
\end{enumerate}
\end{lemma}

\begin{proof}
On a high piece $n_L\ge M/8$, so its rough kinetic term supplies at
least $cM/E^2$.  This pays for the error $Cv_bD^{1-0.001}$ and the
$O(v_b)$ Lee--Huang--Yang tangent correction: the relevant ratio is
$D^{0.001}Q^{-402}\to\infty$.
A fixed fraction of the kinetic term remains.  The scalar quadratic
remainder controls $Dv_b\mathfrak h(n/(Dv_b)-1)$ in the middle range.
It also recenters the microcell penalty at one, since
$\mathfrak h(s+t)\le C(\mathfrak h(s)+\mathfrak h(t))$ and
$\sum_{A\subset b}|A|=v_b$.

On low pieces convexity of $e$ retains the same box penalty in addition
to the two positive forms of Lemma~\ref{lem:energy-sharp}.
Below the middle range the scalar rough bound exceeds $t_b(n)$ by
$cDv_b$, after its error is absorbed.  Above it the grouping argument
used for \eqref{eq:energy-preliminary} gives a surplus $cn$.
These also control the stated box penalty on the outliers.
Box counts commute with the localizations, so their penalties sum back
to the original restrictions.
Finally combine \eqref{eq:energy-upper},
\eqref{eq:energy-mirror-error}, and \eqref{eq:energy-cutoff-error}.
The total scalar errors are
$O(VQ^{-1/4}+VQ^{-1}+VD^{-c_*})$, which are
$O(VQ^{-1/5})$ after making $u$ small relative to $c_*$.
\end{proof}

At $E=R$, the first conclusion immediately gives a clipped second
moment of $q_b-1$ bounded by $CD^{-1}Q^{-1/5}$.  We next prove the
microcell fourth moment at $E=B$.  This is where retaining the
excitation gap, rather than only the scalar energy, is useful.

For a microcell $A\subset b$, put
\[
 X_A=n_A-n|A|/v_b
     =\dGamma\bigl(\1_A-|A|/v_b\bigr).
\]
Its one-body summand has norm at most one and zero constant-mode
diagonal.  Hence for every $n$-particle vector $\eta$,
\begin{equation}\label{eq:energy-fourth-block}
                    \|X_A^2\eta\|^2
          \le Cn^2\langle(n_++2)^2\rangle_\eta.
\end{equation}
Indeed its three blocks relative to $P\oplus P^\perp$ change $n_+$
by $1,-1,0$.  Between the $k$ and $k+1$ sectors their norms are bounded
by $\sqrt{(n-k)(k+1)}$, and the diagonal norm is at most $k$.
Thus $X_A^2$ has bandwidth two and blocks of norm at most $Cn(k+2)$.
Squaring and using Cauchy--Schwarz over the bounded overlaps proves
\eqref{eq:energy-fourth-block}.

Lemma~\ref{lem:energy-surplus} gives
\begin{equation}\label{eq:energy-gap-counts}
 \frac1V\sum_{\rm low}\langle n_+\rangle\le CQ^{2-1/5},\qquad
 \frac1V\sum_{\rm low}\langle n_H\rangle\le CQ^{2-1/5}/J.
\end{equation}
Combine these with \eqref{eq:energy-number-square} and
\eqref{eq:energy-fourth-block}.  Dividing by $(D|A|)^4$ and uniformly
averaging all microcells gives the bound
\begin{equation}\label{eq:energy-microcell-error}
 C_j\left(D^{-1}Q^{-386-1/5}+D^{-1}Q^{-86-1/5}
                         +D^{-2}Q^6+D^{-2}Q^{11}\right)
\end{equation}
for the low-piece fourth moment centered at the relative box mean.
To check the volume factors, $n^2\le CD^2Q^6$, and averaging the
$O_j(Q^3)$ microcells of each box gives prefactor at most
$C_jD^{-2}Q^9/V$ before the summed excitation moments.
The terms $Mn_+$ and $nn_H$ then give the first two powers in
\eqref{eq:energy-microcell-error}; constants and the linear term in
$(n_++2)^2$ give the last two.

Fix a finite clipping level, allowed to depend on $j$.
On the middle number range, a clipped fourth power of the box mean
minus one is bounded by a constant times its square, so the box penalty
recenters the preceding estimate at one.  On high pieces a clipped
fourth power of $q_A-1$ is bounded by $C_j\mathfrak h(q_A-1)$.
On outlying number sectors it is bounded by a constant times the norm
penalty of Lemma~\ref{lem:energy-surplus} after averaging cells.
Since $\psi=F_1\psi+F_2\psi$, the triangle inequality applied to the
square root of the clipped nonnegative observable bounds its
expectation in the original middle vector by twice the sum of those
in $F_1\psi$ and $F_2\psi$.  The extra low piece $F_3\psi$ needed for
energy localization has a nonnegative contribution and may be omitted
here.  This proves the clipped fourth-moment bound for the original
ground state.  The estimates beat $D^{-1}Q^{-1/6}$ for sufficiently
small $u$.

We remove clipping only at this point, in the original state where
\eqref{eq:count-tails} applies.  Choose a large fixed unit-count threshold
$C_0D$.  Except on the event that some constituent unit cell exceeds
that threshold, an $R$-cell has bounded relative count and a $d$-cell
has relative count at most $C_0d^{-3}$.  A fixed count moment is
polynomial in $D$, also for regions with polynomially many unit cells,
by \eqref{eq:count-mgf}.  Cauchy--Schwarz and the exponential upper tail
therefore bound every discarded second or fourth moment by a quantity
smaller than every inverse power of $D$.  This proves both count
assertions in \eqref{eq:local-energy}, with their stated slightly weaker
exponents.

\subsection{Recovering the occupation from the quasiparticle vacuum}

It remains to extract the two occupation statements.  Use $E=B$ in
this subsection.  For $l=R$ or $d$, the one-body operator $W_l$ acts
within each $B$-box, annihilates its constant, and satisfies
\begin{equation}\label{eq:energy-projector-bounds}
 0\le W_l\le1,\qquad W_l\le Cl^2(-\Delta_{\rm Neu}).
\end{equation}
The second inequality is the Poincar\'e inequality on its component
$l$-cells.  On high pieces these bounds control $\dGamma(W_l)$ by
$C(1+R^2)\dGamma(\min\{p^2,1\})$: split at $p^2=1$ and use form
Cauchy--Schwarz for the cross term.  Lemma~\ref{lem:energy-surplus}
then makes the contribution $o(V)$ because $R^2B^{-1/5}=o(1)$.
Outlying sectors contribute $o(V)$ by their norm penalty.

For a low piece at a coherent parameter, the Bogoliubov coefficients
satisfy
\begin{equation}\label{eq:energy-uv-control}
 |V_p|^2\le C\min\{|p|^{-1},|p|^{-4}\},\qquad
 p^2|U_p|^2\le CD_p\quad(p\text{ low}).
\end{equation}
For high $p$, $U_p$ is bounded and
$\widetilde D_p\ge cJ/E^2$.  Define the vacuum trace and positive form
\[
 T_l(z)=\sum_{p\ne0}(W_l)_{pp}|V_p|^2,\qquad
 \mathcal Q_z(\eta)=\sum_{p\ne0}\widetilde D_p\|b_p\eta\|^2.
\]
Shell summation gives $T_l(z)\le Cv_b$, uniformly in the allowed $z$.
Substituting \eqref{eq:energy-bogoliubov} into the observable yields
\begin{equation}\label{eq:energy-observable-comparison}
\begin{split}
 \left|\langle\dGamma(W_l)\rangle_\eta-T_l(z)\|\eta\|^2\right|
 \le{}&C(1+R^2)\mathcal Q_z(\eta)\\
 &+C\bigl[(1+R^2)\mathcal Q_z(\eta)T_l(z)\|\eta\|^2\bigr]^{1/2}.
\end{split}
\end{equation}
Here is a useful way to see all terms.  Apply $W_l^{1/2}$ to the
column $(U_pb_p\eta)_p$.  Its squared norm is bounded by the first
term on the right, using \eqref{eq:energy-projector-bounds} and
\eqref{eq:energy-uv-control} on the low band, and $W_l\le1$ on the
high band.  For the column $(V_pb_p^*\eta)_p$, commute $b$ past $b^*$.
The commutator is exactly $T_l(z)\|\eta\|^2$, and the remaining
excitation term has the same bound, also for the transposed matrix.
Cauchy--Schwarz bounds the two cross terms and gives
\eqref{eq:energy-observable-comparison}.

Integrate \eqref{eq:energy-observable-comparison} over coherent
parameters and sum low pieces.  By Lemma~\ref{lem:energy-surplus} its
right side is $o(V)$; for the square-root term use Cauchy--Schwarz and
$\sum v_b\|\psi\|^2\le V$.  Coherent resolution is exact for this
observable, and its large-$z$ tail is negligible by the particle-number
bound.  The one-body localization error already controlled after
\eqref{eq:energy-cutoff-error} is also negligible.  We have therefore
reduced both occupations to the averaged traces $T_l(z)$.

To evaluate them, weight each low coherent piece on $|z|^2\le2n$ by
$v_b\|\psi_z\|^2/V$ and include its spectator and sector integrals.
The resulting finite measure has mass tending to one, and
\begin{equation}\label{eq:energy-coherent-concentration}
                         \rho_z/D\longrightarrow1
             \quad\text{in its measure}.
\end{equation}
Indeed the retained gain $cM/E^2$ bounds the high-piece mass by
$CD^{-1}Q^{402-1/5}=o(1)$; the number penalty controls outlier mass,
and the omitted coherent tail is exponentially small.  The box penalty makes
$n/(Dv_b)$ concentrate at one; \eqref{eq:energy-gap-counts} makes
$n_+/(Dv_b)$ negligible.  Finally, after angular integration the
coherent radial variables have gamma shapes $n-n_++1$ and variances
of that same order.  Their fluctuations divided by $Dv_b$ vanish.
This proves \eqref{eq:energy-coherent-concentration} with the prescribed
order of limits.

On every compact momentum annulus, shrinking support of $g$,
$\tau_p\to p^2$, and \eqref{eq:energy-coherent-concentration} show that
$|V_p|^2$ converges in this average to
\[
 F_\alpha(p)=\frac12\left(
       \frac{p^2+8\pi\alpha}{\sqrt{p^4+16\pi\alpha p^2}}-1\right).
\]
The summable majorant in \eqref{eq:energy-uv-control} removes the small
and large momentum tails, uniformly in the coherent parameters.
Neumann Riemann sums then converge with full-space measure
$\dd\bar p=(2\pi)^{-3}\dd p$.

For $W_R$, its diagonal entries tend uniformly to one on compact
annuli.  In fact the average of a cosine over a cell of side $R$ tends
to zero when one of its frequencies is bounded away from zero; on an
annulus at least one component has this property, and the product
cosine formula gives the uniform assertion.  For $W_d$,
\eqref{eq:energy-projector-bounds} gives
$(W_d)_{pp}\le\min\{1,Cd^2p^2\}$.  Thus its limiting trace density is
bounded by
\[
 \int_{\R^3}\min\{1,Cd^2p^2\}
               C\min\{|p|^{-1},|p|^{-4}\}\,\dd\bar p,
\]
which tends to zero as $d=1/j\to0$ by dominated convergence.
For $W_R$ the full limiting integral is
\begin{align*}
 \int_{\R^3}F_\alpha(p)\,\dd\bar p
 &=\frac{(8\pi\alpha)^{3/2}}{4\pi^2}
   \int_0^\infty\left(
          \frac{x(x^2+1)}{\sqrt{x^2+2}}-x^2\right)\dd x\\
 &=\frac8{3\sqrt\pi}\alpha^{3/2}.
\end{align*}
The radial integral is $\sqrt2/3$, again by elementary integration.
This completes the proof of Proposition~\ref{prop:local-energy}.
